\documentclass[10pt]{article}

\usepackage[utf8]{inputenc}

\usepackage[T1]{fontenc}

\usepackage{amsmath}

\usepackage{amsfonts}

\usepackage{amssymb}

\usepackage[version=4]{mhchem}

\usepackage{stmaryrd}

\usepackage{bbold}

\usepackage{graphicx}

\usepackage[export]{adjustbox}

\graphicspath{ {./images/} }



\begin{document}

Так, например:



\begin{enumerate}

  \setcounter{enumi}{2}

  \item если $p$ - нечетное простое число, то $p$-примарная компонента группы $\pi_{k}^{\mathrm{s}}$ есть $\mathbb{Z}_{p}$ при $k=2 l(p-1)-1$, $l=1, \ldots,(p-1)$, и тривиальна при других $k<2 \mathrm{p}(p-$ $-1)-2$.

\end{enumerate}



Имеется большое количество результатов о С. г. г., область действия к-рых не ограничивается никаким конечным диапазоном значений $i-n$. В частности, известно большое количество бесконечных серий нетривиальных әлементов групп $\pi_{i}\left(S^{n}\right)$ (см. [4]).



\begin{enumerate}

  \setcounter{enumi}{3}

  \item Порядок образа Уайтхеда гомоморфизма $J_{k} \mathrm{pa}-$ вен знаменателю несократимой дроби, равной $B_{k} / 4 k$, где $B_{k}$ есть $k$-е Бернулли число. В частности, Card $\operatorname{Im} J_{1}=24$, Card $\operatorname{Im} J_{2}=240$, Card Im $J_{3}=504$, Card $\operatorname{Im} J_{4}=480$.

\end{enumerate}



Лuт.: [1] П о н т р я г и н Л. С., Гладкие многообразия и их применения в теории гомотопий, 2 изд., М., 1976; [2]

A d a m s J., "Тороlogy», 1963, v. 2, p. 181-95; 1965, v. 3,



\includegraphics[max width=\textwidth, center]{2023_07_09_2729497d7544d34dad64g-1}

Композиционные методы в теории гомотопических групп сфер, пер. с англ., М., 1982; [4] У а й т х е д Д ж., Новейшие достижения в теории гомотопий, пер. с англ., м., 1974; [5] Фомотопическая топология, 2 изд., м., 1969.



СФЕРА - множество $S^{n}$ точек $x$ евклидова пространства $E^{n+1}$, находящихся от нек-рой точки $x_{0}$ (г е н т р С.) на постоянном расстоянии $R$ (р а д иy c C.), T. e.



\[

S^{n}=\left\{x \in E^{n+1}: \rho\left(x, x_{0}\right)=R\right\}

\]



C. $S^{0}-$ пара точек, С. $S^{1}-$ это окружность, C. $S^{n}$ при $n>2$ иногда наз. г и п е р с ф е р о й. Объем С. $S^{n}$ (длина при $n=1$, поверхность при $n=2$ ) вычисляется по формуле



в тастности,



\[

S^{n}=\frac{2 \pi^{\frac{n+1}{2}}}{\Gamma\left(\frac{n+1}{2}\right)} R^{n},

\]



$S^{1}=2 \pi R, S^{2}=4 \pi R^{2}, S^{3}=2 \pi^{2} R^{3}, S^{4}=\frac{8}{3} \pi^{2} R^{4}$.



У равнение С. $S^{n}$ в декартовых прямоугольных координатах в $E^{n+1}$ имеет вид



\[

\sum\left(x^{i}-x_{0}^{i}\right)^{2}=R^{2}

\]



(здесь $x^{i}, x_{0}^{i}, i=1, \ldots, n+1,-$ координаты $x, x_{0}$ соответственно), т. е. С. - (гипер)квадрика, или поверхность второго порядка специального вида.



Положение какой-либо точки в пространстве относительно С. характеризуется степенью точки. Совокупность всех С., относительно к-рых данная точка имеет одинаковую степень, составляет сеть С. Совокупность всех С., относительно к-рых точки нек-рой прямой (р а д и к а л ь н о й о с и) имеют одинаковую степень (различную для различных точек), составляет связку С. Совокушность всех С., относительно к-рых точки нек-рой плоскости (р а д и к а ль н о й п л о с к ос т и) имеют одинаковую степень (различную для разных точек), составляет пучок $\mathrm{C}$.



C точки зрения дифференциальной геометрии, С. $S^{n}$ - риманово пространство, имеющее постоянную (гауссову при $n=2$ и риманову при $n>2$ ) кривизну $k=\frac{1}{R^{n}}$. Все геодезич. линии С. замкнуты и имеют постоянную длину $2 \pi R$ - это т. н. боль ши е окр у ж н о с ти, т. е. пересечения с $S^{n}$ двумерных плоскостей в $E^{n+1}$, проходящих через ее центр. Внешнегеометрич. свойства $S^{n}$ : все нормали пересекаются в одной точке, кривизна любого нормального сечения одна и та же и не зависит от точки, в к-рой оно рассматривается, в частности имеет постоянную среднюю кривизну, причем полная средняя кривизна С. - наименьшая среди выпуклых поверхностей одинаковой площади, все точки С. омбилические. Нек-рые из таких свойств, принятые за основные, послужили отправной точкой для обобщения понятия C. Так, напр., аффинная сфера определяется тем, что все ее (аффинные) нормали пересекаются в одной точке; псевдосфера - поверхность в $E^{3}$ постоянной гауссовой кривизны (но уже отрицательной); одна из интерпретаций орисферы (п р е д е л ь н о й с ф е р ы) - множество точек внутри $S^{2}$, определяемое уравнением также второго порядка



\[

\left(1-x^{2}-y^{2}-z^{2}\right)=\operatorname{const}(1-x \alpha-y \beta-z \gamma)^{2} .

\]



На C. $S^{n}$ дважды транзитивно действует ортогональная группа $O(n+1)$ пространства $E^{n+1}$ (2-транзитивность означает, что для любых двух пар точек с равными расстояниями между ними существует вращение - элемент $O(n+1)$, переводящее одну пару в другую); эта группа является полной группой изометрий $S^{n}$; наконец, С. есть однородное пространство: $S^{n}=$ $=O(n+1) / O(n)$.



$\mathrm{C}$ точки зрения (дифференциальной) топологии, С. $S^{n}$ - замкнутое дифференцируемое многообразие, разделяющее $E^{n+1}$ на две области и являющееся их общеӥ границей; при этом ограниченная область, гомеоморфная $E^{n+1}$,- это (открытый) uар, так что С. можно определить как его границу.



Групшы гомологий С. $S^{n}, n \geqslant 1$ :



\[

H_{k}\left(S^{n}\right)= \begin{cases}0, & k \neq n, \\ \mathbb{Z}, & k=n,\end{cases}

\]



в частности $S^{n}$ не стягивается в точку сама по себе, т. е. тождественное отображение $S^{n}$ в себя существенно.



Группы гомотопий С. $S^{n}, n \geqslant 1$ :



\[

\pi_{k}\left(S^{n}\right)= \begin{cases}0, & k<n, \\ \mathbb{Z}, & k=n .\end{cases}

\]



Напр., $\pi_{3}\left(S^{2}\right)=\mathbb{Z}, \pi_{n+1}\left(S^{n}\right)=\pi_{n+2}\left(S^{n}\right)=\mathbb{Z}_{2}$ при $n>2$. В общем случае - для любых $k$ и $n, k>n$, группы $\pi_{k}\left(S^{n}\right)$ не вычислены (см. Сфер гомотопические группы).



И здесь понятие С. получает обобщение. Напр., дикая сфера - топологич. С. (см. ниже) в $E^{n+1}$, не огравичивающая области, гомеоморфной $E^{n+1} ;$ Милноpa cøepa (э к з о т и ч е с к а я С.) - многообразие, гомеоморфное, но не диффеоморфное $S^{n}$.



Топологич. пространство, гомеоморфное С., наз. то поло ги ч е с к о й сфе р о й. Одним из основных здесь является вопрос об условиях того, что нек-рое пространство является топологич. сферой.



П р и м е р ы. а) Инвариантная топологич. характеристика С. $S^{n}$ при $n>2$ не известна (1984). О случае $n=1$ см. Одномерное многообразие. Для төго ттобы континуум был гомеоморфен C. $S^{2}$, необходимо и достаточно, чтобы он был локально связан, содержал хотя бы одну простую замкнутую линию и чтобы всякая лежащая на нем такая линия разбивала его на две области, имеющие эту линию своей общей границей (т е о р ем а У а й л д е p a).



б) Полное односвязное риманово пространство размерности $n \geqslant 2$, кривизна $K_{\delta}$ к-рого для всех касательных двумерных плоскостей $\sigma \delta$-ограничена с $\delta>1 / 4$, т. е. $\delta \leqslant K_{\delta} \leqslant 1$ гомеоморфно $S^{n} \quad$ (т е о р е м а о с ф ер е, см. Риманова геометрия).



в) Односвязное замкнутое гладкое многообразие, (делые) гомологии к-рого совпадают с гомологиями $S^{n}$, гомеоморфно $S^{n}$ при $n \geqslant 4$ (при $n=3$-неизвестно (1984)). Если $n=5,6$, то оно также и гомеоморфно $S^{n}$ (о б о бшенная ги по те з а П у ан к а р е), при $n=3,4$ гипотеза остается (1984), при $n \geqslant 7$ диффеоморфизм не имеет места.



Совершенно аналогично определяется С. $S$ в метрич. пространстве $(M, \rho): S=\left\{x \in M, \rho\left(x, x_{0}\right)=R\right\}$.





\end{document}